\subsection{A uniform weighted interaction bound}
\label{quantum-clusters:effective}

Let \(V=[n]\) index finite-dimensional sites.  Consider a Hermitian
Hamiltonian
\[
 H=\sum_{A\in\mathcal A}h_A,\qquad
 1\le |A|\le k,\qquad
 \Lambda=\max_{i\in V}\sum_{A\ni i}\|h_A\|,
 \tag{QH}
\]
where \(h_A\) is supported on \(A\), \(\mathcal A\) is finite, and the
norm is the operator norm.  Multiple labelled terms with the same
support are allowed.  Scalar terms may be discarded.  No commutativity,
geometric locality, or degree bound is assumed.  In particular,
\(\Lambda\) includes one-site terms.

Fix \(\beta\ge0\).  For each site and each \(x_i\in\{-1,1\}\), let \(\varphi_{i,x_i}\)
be a positive normalized linear functional on the local operator
algebra, and put
\[
 \varphi_x=\bigotimes_{i\in V}\varphi_{i,x_i},\qquad
 F_\beta(x)=\varphi_x(e^{-\beta H}),\qquad
 V_\beta(x)=\log F_\beta(x).
\]
Positivity and normalization imply \(|\varphi_x(B)|\le\|B\|\),
and \(F_\beta(x)>0\) for real \(\beta\).  Define Walsh coefficients
using the uniform measure on \(\{-1,1\}^V\).

\begin{theorem}[Noncommuting interactions give a small classical potential]
\label{thm:quantum-clusters-potential}
Assume $\beta\ge0$, and set
\[
 c_* = 2+\tfrac12\log 2,
 \qquad L=\beta\Lambda e^{c_* k}.
\]
If \(e k L<1\) and
\(R:=eL/(1-ekL)\le1\), then
\[
 \max_i\sum_{S\ni i}(|S|-1)
       |\widehat V_\beta(S)|\le R.
 \tag{QW}
\]
In particular,
\[
 \boxed{\quad
 \beta\Lambda\le
 \frac{e^{-k(2+\frac12\log2)}}{32ek}
 \quad\Longrightarrow\quad
 \max_i\sum_{S\ni i}(|S|-1)
       |\widehat V_\beta(S)|\le\frac1{31k}<\frac18.
 \quad}
 \tag{QT}
\]
All bounds are independent of \(n\), the dimensions of the local
Hilbert spaces, the number of terms incident to a site, and the chosen
product functionals.  The constants are sufficient bounds, without a
claim of optimality.
\end{theorem}

\subsubsection{The exact polymer representation}

A polymer \(\gamma\) is a nonempty multiset of term labels whose
occurrence-overlap graph is connected.  Its multiplicities are
\(m_A\), its total multiplicity is \(m=\sum_A m_A\), and its support
is \(B_\gamma=\bigcup_{m_A>0}A\).  The occurrence-overlap graph has
one vertex per occurrence; two vertices are joined when their term
supports intersect.  Two polymers are incompatible exactly when their
supports intersect, including self-incompatibility.

Let \(\mathcal O(\gamma)\) be the set of distinct orderings of its
multiset of term labels.  Define the function
\[
 w_\gamma(x)=\frac{(-\beta)^m}{m!}
    \sum_{(A_1,\ldots,A_m)\in\mathcal O(\gamma)}
        \varphi_x(h_{A_1}\cdots h_{A_m}).
 \tag{QA}
\]
It depends only on \(x_{B_\gamma}\).  Because
\(|\mathcal O(\gamma)|=m!/\prod_A m_A!\),
\[
 a_\gamma:=\|w_\gamma\|_\infty
 \le\prod_A\frac{(\beta\|h_A\|)^{m_A}}{m_A!}.
 \tag{QM}
\]
The function is real: reversal pairs each word with its adjoint, while
the sum includes both words.  The arguments below also allow complex
activities, so reality is not required for convergence.

Expanding \(e^{-\beta H}\) in ordered words and separating each word
multiset into its connected overlap components gives the exact identity
\[
 F_\beta(x)=
 \sum_{\substack{\Gamma\text{ a finite family of polymers}\\
                  B_\gamma\cap B_{\gamma'}=\varnothing
                  \text{ for }\gamma\ne\gamma'}}
                  \prod_{\gamma\in\Gamma}w_\gamma(x).
 \tag{QP}
\]
Here the empty family contributes one.  To verify the coefficient,
fix internal ordered words of lengths \(m_1,\ldots,m_r\) in disjoint
components.  There are \(m!/(m_1!\cdots m_r!)\) interleavings.
Operators on the distinct supports commute, and \(\varphi_x\) factors
over them.  Thus the Taylor denominator \(m!\) becomes exactly the
product of denominators in (QA).  Every multiset has a unique component
decomposition, so there is no additional multiplicity.  Absolute
convergence follows already from the Taylor majorant
\(\exp(\beta\sum_A\|h_A\|)\) for fixed finite \(V\).

\subsubsection{A connected-polymer count without a degree bound}

For \(c\ge0\), put
\[
 R_i(c)=\sum_{\gamma:\,i\in B_\gamma}a_\gamma e^{c|B_\gamma|},
 \qquad L_c=\beta\Lambda e^{ck}.
\]
Then
\[
 R_i(c)\le
 L_c\sum_{m\ge1}\frac{m^{m-1}}{m!}(kL_c)^{m-1}
 \le\frac{eL_c}{1-ekL_c}
 \quad\text{if }ekL_c<1.
 \tag{QR}
\]
We give the count to make explicit why no unweighted degree enters.
Write \(t_A=\beta\|h_A\|e^{c|A|}\).  Since
\(|B_\gamma|\le\sum_A m_A|A|\), (QM) is majorized after weighting
by \(\prod_A t_A^{m_A}/m_A!\).  Multisets of size \(m\) can be
summed as \(1/m!\) times a sum over ordered \(m\)-tuples of labels,
including repetitions.  A support containing \(i\) has at least one
occurrence containing \(i\), so root the occurrence tree at one of at
most \(m\) choices.  For \(m\ge2\), a connected occurrence graph
has a spanning tree, and there are \(m^{m-2}\) labelled trees.
For a fixed tree, sum label assignments from the leaves towards its
root.  The root-label sum is at most \(L_c\); for each fixed parent
label \(A\),
\[
 \sum_{A':\,A'\cap A\ne\varnothing}t_{A'}
 \le\sum_{j\in A}\sum_{A'\ni j}t_{A'}\le kL_c.
\]
This proves the first bound in (QR); its \(m=1\) term is interpreted
as \(L_c\).  Finally,
\(m^{m-1}/m!\le e^m\) gives the geometric-series bound.

\subsubsection{The cluster expansion and its sitewise bound}

We use the following standard abstract-polymer fact.  If nonnegative
majorants \(z_\gamma\) and weights \(b_\gamma\) obey
\[
 \sum_{\gamma':\,B_{\gamma'}\cap B_\gamma\ne\varnothing}
     z_{\gamma'}e^{b_{\gamma'}}\le b_\gamma,
 \tag{QKP}
\]
the logarithmic polymer expansion converges absolutely and its pinned
derivative satisfies
\[
 \Pi_\gamma(z):=
 \sum_{r\ge1}\frac1{(r-1)!}
    \sum_{\gamma_2,\ldots,\gamma_r}
    |\phi^T(\gamma,\gamma_2,\ldots,\gamma_r)|
       \prod_{j=2}^r z_{\gamma_j}
 \le e^{b_\gamma}.
 \tag{QPIN}
\]
Here \(\phi^T\) is one for \(r=1\), and otherwise is the alternating
sum over connected spanning subgraphs of the incompatibility graph.
This is the Koteck\'y--Preiss bound in the pinned-derivative formulation;
see Fern\'andez--Procacci,
\href{https://arxiv.org/pdf/math-ph/0605041}{equations (2.7) and
(2.12)--(2.15)}.  Their result permits countably many polymers,
self-incompatibility, and complex activities dominated by \(z\).

Let \(t_*=1+\frac12\log2\), so \(c_*=t_*+1\).  Apply this
criterion with
\[
 z_\gamma=a_\gamma e^{t_*|B_\gamma|},
 \qquad b_\gamma=|B_\gamma|.
\]
By (QR) and the theorem's assumptions,
\[
 \sum_{\gamma'\not\sim\gamma}z_{\gamma'}e^{b_{\gamma'}}
 \le\sum_{i\in B_\gamma}R_i(c_*)
 \le |B_\gamma|R\le b_\gamma.
\]
Thus (QKP) holds.  For an ordered cluster
\(C=(\gamma_1,\ldots,\gamma_r)\), write
\(B_C=\bigcup_{j=1}^r B_{\gamma_j}\).  Its contribution to
\(\log F_\beta\) is
\[
 f_C(x)=\frac{\phi^T(\gamma_1,\ldots,\gamma_r)}{r!}
             \prod_{j=1}^r w_{\gamma_j}(x).
 \tag{QC}
\]
Summing over all ordered clusters, including repeated polymers, is the
usual convention.  Because \(|B_C|\le\sum_j|B_{\gamma_j}|\),
rooting at any member whose support contains \(i\) and using (QPIN)
gives
\begin{align*}
 \sum_{C:\,i\in B_C}e^{t_*|B_C|}\|f_C\|_\infty
 &\le \sum_{\gamma:\,i\in B_\gamma}z_\gamma\Pi_\gamma(z)\\
 &\le \sum_{\gamma:\,i\in B_\gamma}
               a_\gamma e^{c_*|B_\gamma|}
 \le R.
 \tag{QS}
\end{align*}
For clarity, the first inequality may count a cluster more than once.
Summing the possible root positions changes \(1/r!\) into
\(1/(r-1)!\), which is exactly the pinned series, so the direction
of the bound is correct.

These estimates justify the identity
\(V_\beta=\sum_C f_C\) with uniform absolute convergence on the
finite cube.  Indeed, (QS) summed over sites bounds the total absolute
series by \(nR\).  The logarithm is the branch equal to zero at
\(\beta=0\); since \(F_\beta>0\) on the real interval from zero
to the stated \(\beta\), it is the ordinary real logarithm there.

\subsubsection{From local cluster functions to Fourier interactions}

For any real or complex function \(f\) supported on a set \(B\),
Parseval and Cauchy--Schwarz give
\[
 \sum_{S\subseteq B}|\widehat f(S)|
 \le 2^{|B|/2}\|f\|_{L^2(U_B)}
 \le 2^{|B|/2}\|f\|_\infty.
\]
Therefore, for \(i\in B\),
\[
 \sum_{S\ni i}(|S|-1)|\widehat f(S)|
 \le (|B|-1)2^{|B|/2}\|f\|_\infty
 \le e^{t_*|B|}\|f\|_\infty.
\]
The last bound is harmless when \(|B|=1\), where the left side is
zero.  Apply this to each \(f_C\), sum absolute values, and invoke
(QS).  This proves (QW).  Under (QT), \(ekL\le1/32\) and
\(R\le1/(31k)\), proving the final assertion. \qed
